% =====================================================================================================
\section{The sampler}\label{sec:sampler}

This section describes the algorithm as the software (an R package with a C++ core) implements it in its default
configuration, and how to read what it reports. The user supplies the data, a random seed and a computing allocation;
everything else is a numerical constant or a stated function of $(n,p)$ (Appendix~\ref{app:params} lists every one, with its code location). The last two
subsections say which parts the theorems cover, and describe separately the parallel engine built for $p=10^6$, which no
theorem covers.

\subsection{Why sequential Monte Carlo, and why this start law}\label{sec:intuition}

The posterior is known up to one constant, so the difficulty is to find where its mass is in an enormous, rugged model
space. It is rugged for two reasons: correlated predictors create near-ties between models that swap a true predictor for
its proxy, and a predictor that is useless once the true ones are in can look useful on its own, because it \emph{leaks}
part of their signal. Sequential Monte Carlo starts many models (\emph{particles}) from an easy distribution close to the
target and moves them to the target in small steps, copying the particles the next distribution favours and spreading
them with Markov moves. The \emph{start law} is that easy distribution: it scores every subset of each block of correlated
predictors exactly, so it resolves a true predictor against its proxies, and it can be sampled exactly. What it gets wrong
is leakage across blocks. The local search from its mode undoes that, deleting the spurious predictors one per move; it
uses no random numbers, so it gives a second route to the true model that does not depend on how well the particles
mixed.

\begin{example}[Leakage, with six predictors]\label{ex:leak}
Take $n=100$, true model $\{1,3\}$ with both coefficients $1$, and unit noise. Predictor $2$ is a proxy of $1$
(correlation $0.93$); $4$ is moderately correlated with $3$ ($0.69$); $5$ is correlated with $1$ ($0.63$); and $6$ is a
proxy of $5$ ($0.89$). With blocks $\{1,2\}$, $\{3,4\}$, $\{5,6\}$, the five models with the largest posterior mass are:

\smallskip
\centerline{\small
\begin{tabular}{@{}lccccc@{}}
\toprule
model & $\{1,3,5\}$ & $\{1,3\}$ & $\{1,3,6\}$ & $\{1,3,4\}$ & $\{1,2,3\}$ \\
\midrule
start law $\pi_0$ & $\mathbf{0.517}$ & $0.198$ & $0.067$ & $0.012$ & $0.008$ \\
posterior $\pi(\cdot\mid y)$ & $0.042$ & $\mathbf{0.688}$ & $0.118$ & $0.051$ & $0.037$ \\
\bottomrule
\end{tabular}}
\smallskip

Within its own block, predictor $5$ explains much of $y$ because it carries part of predictor $1$'s signal, and the
start law cannot see that predictor $1$, in another block, already explains it: predictor $5$ leaks, and the start mode is
$\{1,3,5\}$. The proxy $2$, in predictor $1$'s own block, is handled correctly. The posterior puts $0.69$ on $\{1,3\}$, and
the local search from the start mode makes one move, deleting $5$, and stops at $\{1,3\}$, which has no improving
neighbour. The example only illustrates the mechanism. Result~3 shows that, asymptotically and under its assumptions, this
is what happens in general: the start mode contains the true model and possibly some predictors that enter it wrongly
(here predictor~$5$), all members of a small \emph{leakage set} (Section~\ref{sec:a-rm}), and the search deletes them. (Computed with an independent re-implementation of the coded score,
start law and search on this data set; with six predictors the software itself would enumerate all $64$ models.)
\end{example}

\subsection{The algorithm, step by step}\label{sec:algorithm}

A run has seven stages (Figure~\ref{fig:pipeline}). A \emph{step} is one move of the sampler from one temperature to the
next (the temperatures $\xi$ are defined at stages 4 and 5; $\pi_1$ is the posterior).

\begin{figure}[ht]
\centering
\begin{tikzpicture}[node distance=5mm, every node/.style={font=\footnotesize},
  box/.style={draw, rounded corners, align=center, minimum height=10mm, text width=21mm},
  exact/.style={box, fill=blue!12, line width=0.9pt}, stoch/.style={box, fill=orange!18, dashed, line width=0.9pt},
  det/.style={box, fill=green!14, double}]
\node[exact] (pre) {1. standardise;\\ small $\Om$: exact};
\node[exact, right=of pre] (blk) {2. blocks of\\ $\le B$; tables};
\node[exact, right=of blk] (start) {3. start law $\pi_0$};
\node[stoch, right=of start] (cal) {4. calibrate,\\ then freeze};
\node[stoch, below=12mm of cal] (smc) {5. four islands\\ from $\pi_0$ to $\pi_1$};
\node[det, left=of smc] (ls) {6. two local\\ searches};
\node[det, left=of ls] (out) {7. report};
\draw[-{Stealth}] (pre)--(blk); \draw[-{Stealth}] (blk)--(start); \draw[-{Stealth}] (start)--(cal);
\draw[-{Stealth}] (cal)--(smc); \draw[-{Stealth}] (smc)--(ls); \draw[-{Stealth}] (ls)--(out);
\draw[-{Stealth}, dashed, thick] (start.south) -- node[right, pos=0.5]{start mode $\tSz$} (ls.north);
\draw[-{Stealth}, dotted, thick] (pre.south) |- node[left, pos=0.3]{exact case} (out.west);
\end{tikzpicture}
\caption{The stages of a run. Exact computations: solid boxes (1--3). Random: dashed boxes (4--5). Deterministic given
their inputs: double boxes (6--7). When $\Om$ is small the run goes from stage~1 straight to the report (dotted). The
second local search starts from the start mode and does not depend on the particles (dashed arrow).}
\label{fig:pipeline}
\end{figure}

\paragraph{Stage 1: inputs, the exact case and the check before the run.} Invalid inputs stop the run with no model; a constant response, or no retained predictor, stops it with the empty
model and no estimates. If the model space is small, at most $(5\Nc_0+4N_0)\,p$ models with $N_0:=\max(2500,n,p)$ and
$\Nc_0:=\lceil512N_0/2500\rceil$ (Appendix~\ref{app:params}), the method scores every model and computes everything
exactly. Otherwise it estimates the memory the run will need and the worst-case time of the local search from the start
mode, and stops with status \texttt{resource\_limit} if the
estimate exceeds the allocation (8\,GiB and 24 hours by default).

\paragraph{Stage 2: blocks and block tables.} The method partitions the predictors into blocks $C_1,\dots,C_J$ of at most
$B:=\min(\lfloor\sqrt n\rfloor,\lfloor\sqrt p\rfloor,16)$ predictors, by complete-linkage clustering on $1-\abs{\hat r_{jl}}$,
cutting the tree top-down until every cluster fits and packing whole clusters into blocks in leaf order. The partition is a
deterministic function of $X$, and correlated predictors tend to share a block. Within each block it scores every subset
$U$ with $\abs{U}\le K$ and stores the \emph{block table} $h_k(U):=m(U)/m(\emptyset)$, with $h_k(\emptyset)=1$; a block has
at most $2^{16}$ subsets. It also builds a \emph{neighbour graph} of correlated predictors for the swap moves
(Appendix~\ref{app:params}).

\paragraph{Stage 3: the start law and the start mode.} The start law multiplies the size prior by the block tables,
\begin{equation}\label{eq:startlaw}
  \pi_0(S)=\frac{\gamma_0(S)}{\bar Z_0},\qquad \gamma_0(S):=\frac{w(\abs{S})}{\bar w_K}\prod_{k=1}^J h_k(S\cap C_k),\qquad
  \bar Z_0:=\sum_{S\in\Om}\gamma_0(S):
\end{equation}
the posterior one would have if the marginal likelihood, relative to the empty model, split into independent block
contributions. It sees every interaction within a block and none between blocks. A recursion over model sizes, combining
the blocks one at a time, computes $\bar Z_0$ exactly and gives an exact sampler for $\pi_0$ (draw the total size, then
the size of each block, from the last block to the first, then a subset of each block). The \emph{start mode} $\tSz$ is a maximiser of $\pi_0$,
computed by the same recursion with maxima: the smallest size among the maximisers, then the lexicographically smallest
set. The posterior and the start law differ by the \emph{discrepancy}
$\Delta(S):=\log\{m(S)/m(\emptyset)\}-\sum_k\log h_k(S\cap C_k)$, the part of the log marginal likelihood that comes from
interactions between blocks: $\gamma_0(S)e^{\Delta(S)}=e^{F(S)}/m(\emptyset)$.

\paragraph{Stages 4 and 5: the sampler.} The sampler moves through the distributions
\begin{equation}\label{eq:path}
  \pi_\xi\ \propto\ \gamma_0\,e^{\xi\Delta},\qquad 0\le\xi\le1,
\end{equation}
from the start law ($\xi=0$) to the posterior ($\xi=1$), along a grid of temperatures $0=\grid0<\dots<\grid{\Ts}=1$ fixed
by the calibration. An island of $N$ particles starts from $N$ independent draws from $\pi_0$. At each step $t$ it
\begin{enumerate}[label=(\alph*),leftmargin=2em,itemsep=0pt,topsep=2pt]
  \item \emph{reweights} particle $v$ by $\omega^{(t)}_v:=\exp\{(\grid t-\grid{t-1})\Delta(S_v)\}$;
  \item \emph{resamples} the particles multinomially, in proportion to these weights;
  \item \emph{moves} every particle by one round of Markov moves that leave $\pi_{\grid t}$ invariant.
\end{enumerate}
A round is an inclusion update (Gibbs) for each predictor, one uniform swap, $B$ swaps along the neighbour graph, and one
independence move that proposes a fresh draw from $\pi_0$. At $\xi=1$ it runs $\Lbi:=\max(4,\lceil\log_2(n+1)\rceil)$ burn-in rounds
and four measured rounds. The island's inclusion-probability estimate $\PIPhat_{r,j}$ averages, over particles and measured
rounds, the conditional inclusion probability of the Gibbs update rather than the 0/1 indicator; its evidence estimate is
$\Zhat_r:=m(\emptyset)\,\bar Z_0\prod_{t=1}^{\Ts}\{N^{-1}\sum_v\omega^{(t)}_v\}$, the exact normaliser of the start law times
the product of the average reweighting factors.

The \emph{calibration} (stage 4) chooses the grid and the effort before the final run, from one pilot island and a few
rounds of four short diagnostic islands; the grid is chosen so that each step keeps at least $90\%$ effective sample
size, computed on a set of candidate models. The calibration starts at
tier $i=0$ and moves at most once to tier $i=1$, which multiplies all particle numbers by four, when one of its
diagnostics fires (Appendix~\ref{sec:supp-cal} gives the full procedure). Three facts are all that the results use: the calibration's outcome is a function of the data and of its own
random numbers; it is \emph{frozen} before the final run draws any random number; and it \emph{succeeds} when it freezes a
grid, by its main construction or a fallback (otherwise the run stops without estimates). The \emph{final run} (stage 5)
consists of four independent \emph{islands} of $N_i=4^iN_0$ particles at the frozen tier $i$, on the frozen grid, each with
its own random stream.

\paragraph{Stage 6: the local searches.} Two models are \emph{neighbours} if one comes from the other by deleting a
predictor, adding one (while the size is below $K$), or swapping a member for a non-member. A local search scans every
neighbour of the current model and moves to the best one (ties broken lexicographically) if its log score is strictly
higher, for at most a capped number of moves. If a scan finds no strictly better neighbour, the search stops and its end
point is \emph{certified} (a local optimum). If it uses up its move cap first, it scans once more: the end point is
certified if that scan finds no improvement, and uncertified otherwise. Two searches run after the islands: from the best model found so far, with move cap
$\ell_n:=\lceil\log_2(n+p+1)\rceil$; and from the start mode $\tSz$, with move cap $\ell_n+\abs{\tSz}$.

\subsection{What a run reports}\label{sec:outputs-method}

A run ends in one of three ways: \emph{exact} ($\Om$ was enumerated), \emph{completed} (the sampler ran to the end; a
\emph{completed run}), or \emph{stopped} (for example by the resource guard, which stops a run whose time or memory would
exceed its allocation). The first two \emph{produce output}. A stopped run reports no estimates; any partial output is
flagged incomplete, and no result covers it. The method records every model
it scores; the \emph{reported model} $\Shat$ is the best of them by log score (exact ties to the smaller model, then
lexicographically). It is \emph{certified} when the method enumerated $\Om$, or when $\Shat$ ends a local search whose last
scan found no strict improvement. An exact run reports $\Shat$, every $\PIP_j$, the median probability model and $\log Z$.
A completed run pools its islands with equal weights,
\begin{equation}\label{eq:pool}
  \PIPhat_j:=\frac14\sum_{r=1}^4\PIPhat_{r,j},\qquad \Zhat:=\frac14\sum_{r=1}^4\Zhat_r,
\end{equation}
and reports $\Shat$, the $\PIPhat_j$, $\MPMhat:=\{j:\PIPhat_j>1/2\}$, $\log\Zhat$ with the lower bound
$\log(0.05\,\Zhat)$, a flag saying whether $\Shat$ is certified, and the list $\mathcal T$ of the best (at most $40p+1$)
models it scored, with their log scores. Islands are not weighted by their evidence estimates. Diagnostics never change an
output (Appendix~\ref{app:params}).

\paragraph{How to read the outputs.} Each output is an estimate of one posterior summary of
Section~\ref{sec:posterior}, and the theorems of Section~\ref{sec:proved} say when it can be trusted.
\begin{itemize}[leftmargin=1.5em,itemsep=1pt]
  \item \emph{The reported model $\Shat$} is the run's estimate of the highest posterior model. It is the best model the
    run scored, so it equals $\HPM$ when the run has scored $\HPM$; when $\Shat$ is certified, no single addition,
    deletion or swap improves it. Result~3 says when $\Shat=\Sz$; Result~4 says when a certified small $\Shat$ cannot miss
    a true predictor.
  \item \emph{The inclusion probabilities $\PIPhat_j$} estimate $\PIP_j$. Each island averages, over its particles and
    four measured rounds, the probability that the Gibbs update would include $j$ (a Rao--Blackwellised average, less
    noisy than counting inclusions); the four islands are averaged with equal weights. A value near $1$ or $0$ means near-certain
    inclusion or exclusion.
  \item \emph{The median probability model $\MPMhat$} is the set of predictors with $\PIPhat_j>1/2$; it estimates $\MPM$,
    the model that is optimal for prediction under conditions of \citet{BarbieriBerger2004}.
  \item \emph{The evidence $\log\Zhat$} estimates $\log Z$, the marginal probability of the data under the whole prior;
    it is what one compares across priors or data sets. Given the calibration, $\Zhat$ is unbiased for $Z$
    (Result~5), so, by Markov's inequality, the probability that the run completes and still reports a lower bound
    $\log(0.05\,\Zhat)$ above $\log Z$ is at most $0.05$.
  \item \emph{The returned list $\mathcal T$} holds the best models scored, with their exact log scores; renormalising
    the posterior to it gives a second set of estimates (supplement, Corollary~\ref{cor:list}).
\end{itemize}
At moderate $n$ these outputs describe a posterior spread over the true model and many one-predictor extensions of it,
not a point mass: in single synthetic draws at $n=200$ the one-predictor extensions together weighed about as much as the
true model itself (Remark~\ref{rem:moderate-n}; Result~6 explains why).

\subsection{What the theorems cover, and what is engineering}\label{sec:coverage}

The theorems are about the method as coded, but they use only a few properties of each stage. Table~\ref{tab:coverage}
says, stage by stage, what a theorem relies on and what is engineering: a choice that makes the software fast or
robust, that no theorem depends on, and that could be changed without touching any statement.

\begin{table}[ht]
\centering\small
\setlength{\tabcolsep}{4pt}
\caption{Which parts of the method the theorems cover. ``Engineering'' means that no theorem depends on the choice.}\label{tab:coverage}
\begin{tabular}{@{}>{\raggedright\arraybackslash}p{0.2\linewidth}>{\raggedright\arraybackslash}p{0.42\linewidth}>{\raggedright\arraybackslash}p{0.32\linewidth}@{}}
\toprule
Part & What the theorems use & Engineering only \\
\midrule
model, priors, posterior & everything: Results 1, 2 and 6 are about this posterior & --- \\
stage 1 & the exact case (Results 3, 4 and Corollary~\ref{cor:list} include it) & the check before the run and the
resource guard (the theorems assume enough memory and time, Convention~\as{C2p}) \\
stage 2, blocks & blocks of at most $16$ predictors, a deterministic function of $X$ (Result~3) & the clustering rule;
the neighbour graph \\
stage 3, start law & the start law and its mode (Result~3) & --- \\
stage 4, calibration & only that it is frozen before the final run and either succeeds or stops the run; Result~3 holds for
every outcome & every rule inside it: the effective-sample-size target, tiers, triggers, rebuilds, fallback \\
stage 5, islands & each island's evidence estimate (Result~5); inclusion probabilities only on designs without leakage,
the independent ones (supplement, Remark~\ref{thm:outputs}) & particle numbers, burn-in, the mix of moves; Result~3 does not need the
particles to mix \\
stage 6, local searches & the search from the start mode (Result~3) and certification (Result~4), with their move caps & --- \\
outputs & $\Shat$ (Results 3, 4), $\Zhat$ (Result 5), the list (Corollary~\ref{cor:list}) & the diagnostics; the
size of the stores \\
parallel $10^6$ engine & nothing (Section~\ref{sec:supp-1m}) & all of it \\
\bottomrule
\end{tabular}
\end{table}

The design choice that makes Result~3 possible is the local search from the start mode: it uses no random numbers, so the
reported model is right on a data event alone, whatever the particles did. The sampler is then needed for the inclusion
probabilities and the evidence, where the theorems are weaker (Results~5 and the supplement's
Remark~\ref{thm:outputs}).

\subsection{The parallel engine for \texorpdfstring{$p=10^6$}{p=10\^{}6} (engineering; no theorem covers it)}\label{sec:supp-1m}

\emph{Nothing in this subsection is covered by a theorem.} To run at $p=10^6$ the project built a separate parallel engine;
its configuration was frozen on 2026-11-19 (code tag \code{p1m-freeze-F2}, on top of the frozen method
\code{method-v3.1.1-listpip}). It is an engineered variant of the method, not the method: it changes the sampler in the
ways listed in Table~\ref{tab:engine}, and none of the theorems above has been proved for it. Its runs are reported
elsewhere; this edition quotes no result from them.

\paragraph{The frozen configuration.} It runs $R'=16$ islands as separate processes that do not communicate, merged as
the method pools its four. Each island has a fixed $N=2{,}500$ particles ($10^4$ for the runs at $p=10^5$; a planned
$10^6$ fit with $N=10^4$ was never run). The calibration is the pilot alone: no rebuilds ($I_{\max}=0$, the number of
rebuild rounds, set in every run plan; the program's default is $3$), no refinement of the pilot's grid ($m=1$, the
setting \code{refine_m}), and no move to the larger tier. The method's multinomial resampling is kept. Appendix
Table~\ref{tab:engine} lists every change against the method.

\paragraph{How it was checked against the method.} Only at $p=10^4$: ten simulated data sets, two seeds each, run with both
the method and the engine (40 fits), scored against criteria declared before any fit (same best model, same median
probability model, inclusion probabilities and evidence within their Monte Carlo error, predictive laws close, cost).
The engine matched the method on all of them, under the scoring script as corrected before scoring (an independent
review had found a defect in the registered script, which printed a weaker verdict); a re-score with $I_{\max}=0$ also
matched. There was
no run of the method at $p=10^5$ or $10^6$; there the engine was judged only by its own pre-declared precision, cost,
memory and restart criteria. Whether a parallel engine that reproduces the method's output inherits its guarantees, and
what more islands and a fixed particle number do to them, are open (Section~\ref{sec:planned}).
